\section{Discusión}
\subsection{Respuesta al objetivo de reconstrucción}
Los resultados permiten una respuesta condicionada y cuantitativa: las propiedades de las especies de partícula y la abundancia neta observada, junto con una temperatura y un modelo de equilibrio explícitos, permitieron reconstruir por separado bariones y antibariones. La condición inversa incorporó el exceso publicado en el cálculo, en lugar de limitarse a compararlo después con una población simétrica. A 150 MeV se infirieron aproximadamente $1.7222\times10^{37}$ partículas por cm$^3$ de cada lado, con una pequeña diferencia de $4.0110\times10^{29}$ cm$^{-3}$. A 40 MeV, la misma condición neta distinguió $1.9166\times10^{28}$ bariones de $1.4894\times10^{28}$ antibariones por cm$^3$.

La prueba de rango precisó qué información procede del presente. La abundancia neta fija una diferencia; añadir el mismo número de partículas y antipartículas no la modifica. Por ello, las masas y la presencia actual, sin condiciones térmicas adicionales, no determinan cuántos pares existieron. El HRG a temperatura especificada aporta esa relación adicional. La reconstrucción obtenida es una inferencia bajo supuestos comprobables, y no la recuperación única de un registro histórico perdido.

El extremo superior de 150 MeV es el punto inicial elegido para nuestros cálculos dentro del régimen hadrónico considerado. No se identificó una fecha universal en la que comenzara la aniquilación de todas las especies de partícula. Hablar de cantidades \emph{reconstruidas} es preciso; un total de eventos \emph{regenerados} requiere resolver las reacciones inversas en el tiempo.

La amplitud entre las tres listas para la densidad total a 150 MeV fue 1.185\% del valor PDG2025. La sensibilidad a la temperatura de partida fue mayor: cambiar de 150 a 156.5 MeV elevó el total físico reconstruido 54.31\% y su rendimiento 26.95\%. Esto exige declarar la temperatura junto a cada cantidad. La referencia de \citet{Bazavov2019} aporta una escala física para la formación de hadrones; justificar un inicio cinético específico todavía requiere las tasas de reacción y la expansión. Cerca de esa transición, verificar el HRG frente a la ecuación de estado de QCD es necesario para atribuir precisión física al conteo.

\subsection{Qué concuerda con lo publicado}
Nuestro contraste con el límite nucleónico de \citet{Satarov2013} mostró que la aproximación de Maxwell-Boltzmann y la evaluación de Fermi-Dirac difirieron apenas 0.05916\% para protones y neutrones a 150 MeV. La diferencia entre nucleones solos y todo el inventario fue mucho mayor: los otros bariones aportaron 74.42\% del total a esa temperatura, 44.13\% a 100 MeV y 1.525\% a 40 MeV. Estas proporciones calculadas aquí concuerdan con la importancia de las resonancias pesadas a temperaturas altas señalada por esos autores. El contraste identifica qué parte de la diferencia se debe al inventario y qué parte a la aproximación de distribución.

La comparación no reprodujo su evolución cinética ni sus correcciones de interacción. Tampoco ofrece una verificación observacional directa de las densidades tempranas. La literatura aportó una aproximación y un comportamiento de referencia; nuestro código evaluó sus magnitudes bajo los supuestos comunes declarados. A 40 MeV la población total inferida todavía fue 7.97 veces el exceso neto, por lo que el extremo inferior tampoco fue una abundancia residual final.

Las curvas de entropía conservaron la disminución de la referencia reducida y concentraron su mayor discrepancia a 150 MeV. En el límite no relativista, con $x_i=m_i/T\gg1$,
\begin{equation}
\frac{s_i}{T^3}\simeq
\frac{g_i}{(2\pi)^{3/2}}x_i^{3/2}
\left(x_i+\frac{5}{2}\right)e^{-x_i}.
\label{eq:asymptotic}
\end{equation}
Esta supresión explica la pérdida de contribución de las masas altas al enfriar. La referencia reducida carece de inventario completo y código original, por lo que su comparación conserva el alcance de una transcripción. La correlación alta describe la forma común de las series; no establece exactitud física ni valida una hipótesis de origen de la materia.

\subsection{Aporte propio y alcance de la novedad}
El aporte comprobable fue la reconstrucción numérica explícita con inventarios fijados, la conexión inversa con el exceso actual, el conteo separado de partículas y antipartículas, y los controles reproducibles. Las nuevas tablas y proporciones son resultados de nuestros cálculos. Los mecanismos de aniquilación, las distribuciones térmicas y el uso de HRG ya están establecidos en la literatura \citep{Vovchenko2019,Satarov2013}. El límite de información de una diferencia neta es una propiedad matemática del problema.

No se demostró aquí un mecanismo físico nuevo, una nueva especie de partícula ni una predicción independiente de la asimetría o del hidrógeno actual. La comparación realizada permite evaluar concordancias y diferencias cuantitativas, pero no basta para establecer prioridad o novedad frente a toda la literatura. Esa afirmación requeriría una revisión específica de trabajos previos y un resultado que pudiera distinguirse de sus predicciones bajo condiciones equivalentes.

\subsection{Secuencia hasta las abundancias actuales}
La Tabla~\ref{tab:external} declara el papel de cada referencia. La revisión de BBN de 2025 incorporó mediciones de deuterio, incluida la de \citet{Cooke2018}, y recomendó el promedio usado en nuestra conversión. Primero, $\eta$ fijó el rendimiento neto; segundo, el HRG y la temperatura permitieron reconstruir las poblaciones tempranas; tercero, $\eta$ y la temperatura del CMB dieron la densidad bariónica media actual; finalmente, $Y_p$ y D/H permitieron expresar una composición primordial equivalente de H, He y D.

La densidad equivalente de H, $(1.898\pm0.015)\times10^{-7}$ cm$^{-3}$, depende de esas entradas adicionales. Su cálculo no constituye una predicción nuclear independiente del HRG. La coincidencia de la densidad bariónica con la tabla publicada comprueba la conversión con una entrada CMB compartida. Las abundancias nucleares requieren evolución hasta BBN y una red de reacciones \citep{Fields2025}; el censo de hidrógeno neutro actual requiere además procesamiento estelar e ionización.

\begin{table*}[t]
\centering\small
\caption{Referencias publicadas y papel en la reconstrucción. Las incertidumbres corresponden a las fuentes, con redondeo declarado para $T_{\gamma0}$; los límites FIRAS son a 95\% de confianza. Se distingue una entrada del cálculo de un canal de comparación que no se ajustó en este trabajo.}
\label{tab:external}
\setlength{\tabcolsep}{5pt}\renewcommand{\arraystretch}{1.15}
\begin{tabular}{>{\raggedright\arraybackslash}p{.18\textwidth}>{\raggedright\arraybackslash}p{.23\textwidth}>{\raggedright\arraybackslash}p{.22\textwidth}>{\raggedright\arraybackslash}p{.26\textwidth}}\toprule
Observable & Valor publicado & Fuente & Papel en la reconstrucción\\\midrule
$\eta_{10}$ del CMB & $6.12\pm0.04$ & \citet{Fields2025} & Condición neta de la inferencia inversa.\\
Temperatura CMB actual & $2.7255\pm0.0006$ K & \citet{Fixsen2009}, redondeado & Densidad fotónica y normalización actual.\\
Fracción bariónica en $^4$He & $Y_p=0.245\pm0.003$ & \citet{Fields2025} & Conversión aproximada a H y He.\\
Deuterio primordial & $(2.508\pm0.029)\times10^{-5}$ & \citet{Fields2025} & Conversión aproximada a D.\\
Densidad bariónica actual & $(2.515\pm0.017)\times10^{-7}$ cm$^{-3}$ & \citet{Groom2025} & Comprobación de conversión; entrada CMB compartida.\\
Densidad $\Omega_bh^2$ & $0.02237\pm0.00015$ & \citet{Planck2020} & Referencia externa; no se ajustó un modelo cosmológico.\\
Distorsiones espectrales & $|y|<1.5\times10^{-5}$\newline $|\mu|<9\times10^{-5}$ & \citet{Fixsen1996} & Restricción relevante para una futura inyección energética tardía.\\
\bottomrule\end{tabular}
\end{table*}

El factor 7.04 usa la aproximación estándar de entropía por fotón después de la aniquilación $e^+e^-$. Los cálculos precisos de desacoplamiento de neutrinos incluyen correcciones \citep{Froustey2020}; el factor no se trató como un parámetro medido con incertidumbre en este trabajo.

\subsection{Poblaciones de equilibrio e historial de aniquilación}
En una expansión adiabática, $S_{\rm com}=\stot a^3$ se conserva aproximadamente. Un esquema de evolución con aniquilación y producción inversa es \citep{Satarov2013}
\begin{align}
\dot n_B+3Hn_B&=-\mathcal C,\\
\dot n_{\bar B}+3Hn_{\bar B}&=-\mathcal C,\\
\mathcal C&=\langle\sigma v\rangle
\bigl(n_Bn_{\bar B}-n_B^{\rm eq}n_{\bar B}^{\rm eq}\bigr).
\end{align}
Al restar las ecuaciones se obtiene $\dot{\Delta n}_B+3H\Delta n_B=0$. Esto justifica usar el exceso conservado como vínculo retrospectivo bajo los supuestos adoptados. Obtener eventos acumulados exige integrar por separado producción y pérdida sobre tiempo y volumen. Una sucesión de estados de equilibrio no proporciona esos totales.

\subsection{Limitaciones de la reconstrucción}
Se usaron masas nominales, anchos nulos, un gas ideal y una temperatura común, sin ajustes a datos de red ni incertidumbres de masas, espín e interacciones. La condición neta no resolvió la neutralidad de carga y extrañeza simultáneamente. Los controles verificaron precisión numérica, no la exactitud física de esos supuestos. La cercanía de 150 MeV a la transición exige comprobar las correcciones del modelo antes de asignar errores físicos \citep{Bazavov2012}.

Las incertidumbres cosmológicas sólo se propagaron a las conversiones actuales; las poblaciones inversas usaron entradas centrales. La amplitud entre catálogos es sensibilidad a inventarios correlacionados. El estudio no identificó el inicio cinético, no integró tasas de reacción ni generó la asimetría primordial \citep{Sakharov1967}. Dentro de ese alcance, establece cantidades tempranas condicionadas y una forma explícita de contrastarlas con el conocimiento disponible.
